% Part: many-valued-logics
% Chapter: sequent-calculus
% Section: proof-theoretic-notions

\documentclass[../../../include/open-logic-section]{subfiles}

\begin{document}

\olfileid{mvl}{seq}{prf}

\olsection{المفاهيم البرهانية النظرية}

\begin{explain}
كما عرّفنا عددًا من المفاهيم الدلالية المهمة، وهي الصحة واللزوم
وقابلية الإشباع، نعرّف الآن~\emph{مفاهيم برهانية نظرية} مقابلة. ولا
تُعرّف هذه المفاهيم بالاحتكام إلى إشباع~!!{sentence}s في~!!{structure}s،
بل بالاحتكام إلى~!!{derivability} أو~!!{nonderivability} تتابعات معينة.
وكان اكتشاف تطابق هذه المفاهيم اكتشافًا مهمًا. وهذا التطابق هو مضمون
\emph{مبرهنتي السلامة} و\emph{الاكتمال}.
\end{explain}

\begin{defn}[المبرهنات]
تكون~!!^a{sentence}~$!A$ \emph{مبرهنة} إذا وُجد~!!a{derivation} في
$\Log{LK}$ للتتابع~$\quad \Sequent !A$. ونكتب~$\Proves !A$ إذا
كانت~$!A$ مبرهنة، و$\Proves/ !A$ إذا لم تكن كذلك.
\end{defn}

\begin{defn}[!!^{derivability}]
تكون~!!^a{sentence}~$!A$ \emph{!!{derivable} من} مجموعة
!!{sentence}s~$\Gamma$، ونكتب~$\Gamma \Proves !A$، إذا وُجدت مجموعة
جزئية منتهية~$\Gamma_0 \subseteq \Gamma$ ومتتالية~$\Gamma_0'$ من
!!{sentence}s الموجودة في~$\Gamma_0$ بحيث~$\Log{LK}$ !!{derive}s
$\Gamma_0' \Sequent !A$. وإذا لم تكن~$!A$ !!{derivable} من~$\Gamma$،
كتبنا~$\Gamma \Proves/ !A$.
\end{defn}

بسبب قواعد التقليص والإضعاف والتبديل، لا يهم ترتيب~!!{sentence}s في
$\Gamma_0'$ ولا عدد مرات ورودها: فإذا كان التتابع
$\Gamma_0' \Sequent !A$ !!{derivable}، كان
$\Gamma_0'' \Sequent !A$ كذلك لكل~$\Gamma_0''$ تحتوي
!!{sentence}s نفسها التي تحتويها~$\Gamma_0'$. فمثلًا، إذا كان
$\Gamma_0 = \{!B, !C\}$، كانت كل من
$\Gamma_0' = \tuple{!B, !B, !C}$ و
$\Gamma_0'' = \tuple{!C, !C, !B}$ متتالية لا تحتوي إلا
!!{sentence}s الموجودة في~$\Gamma_0$. وإذا كان التتابع الذي يحتوي
إحداهما~!!{derivable}، كان الآخر كذلك، مثلًا:
\begin{prooftree}
  \AxiomC{}
  \Deduce$!B, !B, !C \fCenter !A$
  \RightLabel{\LeftR{\Contraction}}
  \UnaryInf$!B, !C \fCenter !A$
  \RightLabel{\LeftR{\Exchange}}
  \UnaryInf$!C, !B \fCenter !A$
  \RightLabel{\LeftR{\Weakening}}
  \UnaryInf$!C, !C, !B \fCenter !A$
\end{prooftree}
سنقول من الآن فصاعدًا إنه إذا كانت~$\Gamma_0$ مجموعة منتهية من
!!{sentence}s، فإن~$\Gamma_0 \Sequent !A$ أي تتابع يكون سابقه متتالية
من~!!{sentence}s الموجودة في~$\Gamma_0$، وسنضمر تطبيق التقليصات
والتبديلات والإضعافات عند الحاجة.

\begin{defn}[الاتساق]
تكون مجموعة الجمل~$\Gamma$ \emph{غير متسقة} إذا وفقط إذا وُجدت
مجموعة جزئية منتهية~$\Gamma_0 \subseteq \Gamma$ بحيث~$\Log{LK}$
!!{derive}s~$\Gamma_0 \Sequent \quad$. وإذا لم تكن~$\Gamma$ غير
متسقة، أي إذا لم تكن~$\Log{LK}$ !!{derive}~$\Gamma_0 \Sequent \quad$
لكل~$\Gamma_0 \subseteq \Gamma$ منتهية، قلنا إنها~\emph{متسقة}.
\end{defn}

\begin{prop}[الانعكاسية]
\ollabel{prop:reflexivity}
إذا كان~$!A \in \Gamma$، كان~$\Gamma \Proves !A$.
\end{prop}

\begin{proof}
التتابع الابتدائي~$!A \Sequent !A$ !!{derivable}، و
$\{!A\} \subseteq \Gamma$.
\end{proof}

\begin{prop}[الرتابة]
\ollabel{prop:monotonicity}
إذا كان~$\Gamma \subseteq \Delta$ و$\Gamma \Proves !A$، كان
$\Delta \Proves !A$.
\end{prop}

\begin{proof}
افترض~$\Gamma \Proves !A$؛ أي توجد~$\Gamma_0 \subseteq \Gamma$
منتهية بحيث يكون~$\Gamma_0 \Sequent !A$ !!{derivable}. ولأن
$\Gamma \subseteq \Delta$، تكون~$\Gamma_0$ أيضًا مجموعة جزئية منتهية
من~$\Delta$. ومن ثم يبيّن~!!{derivation}~$\Gamma_0 \Sequent !A$ أيضًا
أن~$\Delta \Proves !A$.
\end{proof}

\begin{prop}[التعدي]
\ollabel{prop:transitivity}
إذا كان~$\Gamma \Proves !A$ و$\{!A\} \cup \Delta \Proves !B$، كان
$\Gamma \cup \Delta \Proves !B$.
\end{prop}

\begin{proof}
إذا كان~$\Gamma \Proves !A$، وُجدت~$\Gamma_0 \subseteq \Gamma$
منتهية و!!a{derivation}~$\pi_0$ لـ$\Gamma_0 \Sequent !A$. وإذا كان
$\{!A\} \cup \Delta \Proves !B$، وُجدت مجموعة جزئية منتهية
$\Delta_0 \subseteq \Delta$ و!!a{derivation}~$\pi_1$ لـ
$!A, \Delta_0 \Sequent !B$. تأمل~!!{derivation} الآتي:
\begin{prooftree}
  \AxiomC{}
  \RightLabel{$\pi_0$}
  \Deduce$\Gamma_0 \fCenter !A$
  \AxiomC{}
  \RightLabel{$\pi_1$}
  \Deduce$!A, \Delta_0 \fCenter !B$
  \RightLabel{\Cut}
  \BinaryInf$\Gamma_0, \Delta_0 \fCenter !B$
\end{prooftree}
بما أن~$\Gamma_0 \cup \Delta_0 \subseteq \Gamma \cup \Delta$، يبيّن
ذلك~$\Gamma \cup \Delta \Proves !B$.
\end{proof}

لاحظ أن هذا يعني بوجه خاص أنه إذا كان~$\Gamma \Proves !A$ و
$!A \Proves !B$، كان~$\Gamma \Proves !B$. ويلزم أيضًا أنه إذا كان
$!A_1, \dots, !A_n \Proves !B$ و$\Gamma \Proves !A_i$ لكل~$i$، كان
$\Gamma \Proves !B$.

\begin{prop}
\ollabel{prop:incons}
تكون~$\Gamma$ غير متسقة إذا وفقط إذا كان~$\Gamma \Proves {!A}$ لكل
جملة~$!A$.
\end{prop}

\begin{proof}
تمرين.
\end{proof}

\begin{prob}
أثبت~\olref[mvl][seq][prf]{prop:incons}.
\end{prob}

\begin{prop}[التراص]
  \ollabel{prop:proves-compact}
  \begin{enumerate}
  \item إذا كان~$\Gamma \Proves !A$، وُجدت مجموعة جزئية منتهية
    $\Gamma_0 \subseteq \Gamma$ بحيث~$\Gamma_0 \Proves !A$.
  \item إذا كانت كل مجموعة جزئية منتهية من~$\Gamma$ متسقة، كانت
    $\Gamma$ متسقة.
  \end{enumerate}
\end{prop}

\begin{proof}
  \begin{enumerate}
    \item إذا كان~$\Gamma \Proves !A$، وُجدت مجموعة جزئية منتهية
      $\Gamma_0 \subseteq \Gamma$ بحيث يكون للتتابع
      $\Gamma_0 \Sequent !A$ !!a{derivation}. وبالتالي
      $\Gamma_0 \Proves !A$.
    \item إذا كانت~$\Gamma$ غير متسقة، وُجدت مجموعة جزئية منتهية
      $\Gamma_0 \subseteq \Gamma$ بحيث~$\Log{LK}$ !!{derive}s
      $\Gamma_0 \Sequent \quad$. لكن~$\Gamma_0$ عندئذ مجموعة جزئية
      منتهية غير متسقة من~$\Gamma$.
  \end{enumerate}
\end{proof}

\end{document}
